\section{Introduction}
\label{sec:introduction}

Dissolved oxygen below roughly 6~mg/L can kill farmed sturgeon within hours, not days
\citep{lindholm2023waterquality,agmrc2026waterquality}. A locally hosted large language model is
genuinely useful in this setting -- it can reason jointly over sensor readings and fish-behavior
video, and explain its reasoning in plain language, in a way a fixed threshold table cannot -- but
an LLM that silently misses a safety-critical threshold is not a tool anyone should deploy
unattended. S.U.R.E. is an implemented, real-time-capable welfare-monitoring architecture for
sturgeon in a recirculating aquaculture system (RAS) that answers this tension with a specific
architectural choice: a locally hosted LLM (AQUA-1B, with a LoRA domain adapter) proposes a severity
assessment, but a deterministic rule engine has the final word, and can only ever raise the model's
proposed severity, never lower it. This paper reports what happened when that architecture was
measured, not merely designed -- across the decision layer itself and four other layers of the same
system.

No physical water-quality sensor exists for this system, and none was built for this paper. That
absence shapes what kind of contribution this is, and it is worth stating plainly rather than only
as a caveat: without real sensor hardware, this work demonstrates the architectural feasibility of
running a multi-component pipeline -- a locally hosted LLM, a computer-vision detector, a
retrieval-augmented knowledge base, and a deterministic rule engine -- together, and specifically
tests whether that pipeline's space and memory footprint is viable on realistic edge hardware. Framed
this way, the paper is a feasibility and resource-management study of the architecture, run against
synthetically generated sensor values and real, hand-labeled camera imagery captured from the
system's own physical setup, not a field validation of decision-making accuracy in a live setting.
That is a legitimate and, we think, useful contribution in its own right, distinct from and prior to
the question of whether the model's decisions are accurate against a live water body.

\question{Does escalation-only deterministic authority over an LLM's output survive contact with a
genuinely unreliable local model, in a system that is actually running, rather than only in a
design diagram?}

The proposal this system was originally built around assumed a specific answer: the model proposes
a severity judgment, and the rule engine's value is in catching the cases where that judgment is a
dangerous under-estimate. That is a clean story, and it is not the one this paper's measurements
support.

\question{When the architecture ``works,'' what does that actually mean, once measured --
sophisticated correction of a nuanced misjudgment, or something structurally different?}

Measured directly against all eight of the evaluation harness's built-in scenarios, using the real,
unmodified override logic and the system's configured model, the dominant pathway is not correction. Half the
time, the model's output does not parse into anything usable at all, and the rule engine's entire
contribution is defaulting to a safe state rather than adjudicating a disagreement. Only one of
eight scenarios shows the textbook pattern of the model under-calling a threshold and the rule
engine escalating it. And in a finding this paper considers its most novel single result, two of
the three scenarios where the model's final answer agreed with the rule engine did so via reasoning
that cited a sensor value never present in the input -- fabrication that is invisible to any
evaluation checking only the output label, which is how most guardrail and agentic-benchmark
literature is scored. This was measured, deliberately, against the weakest plausible version of the
model component: the configured adapter trained on eight handwritten examples the project's own
fine-tuning pipeline documents as insufficient, while a properly sized, 128-example alternative sat
unused on disk. That the safety property held regardless is stronger evidence for the
architecture's design than a result obtained against a carefully tuned model would have been, and
this paper states that directly rather than treating the undertrained adapter as a limitation to
apologize for.

\question{Does the same narrow-enumerable-interface discipline generalize across architecturally
distinct layers of one system?}

Four more measurements bear on this question, each earning its place independently rather than
restating the first. Agentic tool-routing was abandoned in favor of deterministic code specifically
because an adversarial re-run -- designed to break the negative result, not confirm it -- only
strengthened it. The vision pipeline's own confidence-threshold choice turned out to run at a
different, and more reassuring, point on its precision--recall curve than the number this field
would otherwise headline. INT8 quantization, widely assumed to be where an edge deployment pays for
its speed, turns out to lose \emph{less} accuracy than the unquantized export formats already do.
And a drift-detection binning choice that credit-scoring literature already documents as sensitive
to convention turns out, on this system's own confidence distribution, to invert the naive
expectation at low severity before confirming it at high severity.

This paper does not claim to have solved LLM unreliability, and that claim would not survive its
own evidence: the LLM component measured here is, on its own terms, frequently unreliable. What the
evidence supports is narrower and, we think, more useful to a practitioner building a similar
system: a falsifiable design rule under which that unreliability becomes harmless rather than
merely mitigated. Architect the pipeline assuming the model will often produce nothing usable,
verify the fallback default is conservative, and never let the validator read anything from the
model beyond the one enumerable field it is designed to certify -- and the model's failure modes,
including the ones this paper newly documents, stop being the system's failure mode.

That discipline is worth taking seriously precisely because this paper's own system still lapsed
once in exactly the way it warns against: during the audit behind this paper, the system's own
retrieval-augmented knowledge base was found to still be serving a stale accuracy figure in one file
it had corrected everywhere else -- a knowledge-base entry the correcting commit missed by 43
minutes and that has sat uncorrected since. That the same audit process caught it is the argument
for the discipline, not a reason to distrust the rest of this paper's numbers -- it is disclosed
here, and again in Section~\ref{sec:known-limitations}, as a live, currently unresolved issue rather
than a historical anecdote.

Two further scope statements belong here, early, alongside the two above, because they bound what
``measured'' means throughout this paper. First, every result reported here is an offline,
bench-level evaluation of the implemented system -- against synthetically generated sensor readings,
static labeled imagery, and scripted or synthetic scenarios run against the codebase -- and not a
field deployment: no experiment in this paper involved real physical sensors or an operating
recirculating aquaculture facility, and only the system's interface, not a live deployment, was ever
demonstrated. Second,
S.U.R.E. was built for the TEKNOFEST Agricultural Technologies competition, but this team did not
enter that competition and did not pass its preliminary evaluation round. Both points are stated
again, in full, in Section~\ref{sec:known-limitations}.

\takeaway{Malformed-output fail-safe defaulting, not sophisticated correction, is the dominant
realized safety pathway (50\% versus 12\% of scenarios) -- Section~\ref{sec:results-dual-layer},
Figure~\ref{fig:exp03-buckets}.}

\takeaway{Fabricated sensor values survive inside outputs whose final status agrees with the rule
engine, invisible to status-only evaluation -- Section~\ref{sec:results-dual-layer}.}

\takeaway{The agentic tool-routing negative result replicates under an adversarial re-run built to
break it, not merely confirm it -- Section~\ref{sec:results-agentic}.}

\takeaway{The vision pipeline's configured operating point governs recall at 0.782, materially
different from the headline academic-argmax recall of 0.719 that the rest of this literature would
otherwise cite -- Section~\ref{sec:results-vision}, Figure~\ref{fig:exp07-pr}.}

\takeaway{INT8 quantization loses less accuracy than fp32 export, relocating the real cost to a
shared post-processing path -- Section~\ref{sec:results-vision}, Figure~\ref{fig:exp06-scatter}.}

The rest of this paper proceeds as follows. Section~\ref{sec:related-work} places this system
against the literature it draws on, pivoting explicitly from the domain and deployment literature
that earns citations but no novelty claim, through the guardrail and safety-critical-LLM literature
that this system's pattern does not invent, to the specific gap -- cross-layer measurement plus a
fabrication-inside-agreement finding -- that this paper actually fills.
Section~\ref{sec:background} defines the terms used throughout. Section~\ref{sec:system-architecture}
describes the dual-layer decision system, the vision, retrieval, and MLOps layers, and states the
four formal properties the rest of the paper depends on.
Section~\ref{sec:experimental-setup} documents how every one of the eleven underlying experiments
was actually run. Section~\ref{sec:results} is the paper's evidentiary core.
Section~\ref{sec:discussion} and Section~\ref{sec:known-limitations} state what this paper does and
does not license a reader to conclude, and Section~\ref{sec:conclusion} closes with the paper's
central claim restated against everything measured in between.
